\documentclass[12pt,a4paper]{article}
\usepackage[margin=2cm]{geometry}
\usepackage[T2A]{fontenc}
\usepackage[utf8]{inputenc}
\usepackage[russian]{babel}
\usepackage{amsmath,amssymb,mathtools}
\usepackage{array,booktabs,longtable}
\usepackage{xcolor}
\usepackage{hyperref}
\usepackage{tikz}

\hypersetup{
  colorlinks=true,
  linkcolor=blue!55!black,
  urlcolor=blue!55!black,
  pdftitle={Ревью домашней работы},
  pdfauthor={Лёвин Артём Александрович}
}

\setlength{\parindent}{0pt}
\setlength{\parskip}{0.7em}
\renewcommand{\arraystretch}{1.25}

\begin{document}

\begin{titlepage}
\centering
\vspace*{0.28\textheight}
{\LARGE\bfseries Ревью домашней работы\par}
\vspace{2.2cm}
{\large Автор: Лёвин Артём Александрович\par}
\vspace{0.8cm}
{\large 16 сентября 2026 года\par}
\vfill
\begin{tikzpicture}
  \draw[blue!40!black, line width=0.6pt]
    (-3,0) .. controls (-1.5,0.35) and (1.5,-0.35) .. (3,0);
\end{tikzpicture}
\end{titlepage}

\newpage
\section*{Сводная таблица}
\small
\begin{tabular}{>{\centering\arraybackslash}p{0.07\linewidth} p{0.18\linewidth} p{0.31\linewidth} p{0.13\linewidth} p{0.13\linewidth}}
\toprule
№ задания & Тема & Суть ошибки & Ваш ответ & Правильный ответ \\
\midrule
\hyperref[task:7]{7} & Степени с одинаковым основанием & При умножении степеней с одинаковым основанием показатели должны складываться. В решении сумма показателей $-2+2$ не была получена как ноль. & $0{,}6$ & $1$ \\
\bottomrule
\end{tabular}
\normalsize

\newpage
\vspace*{\fill}
\begin{center}
{\LARGE\bfseries Разбор ошибок}
\end{center}
\vspace*{\fill}

\newpage
\phantomsection\label{task:7}
\section*{Задание №7}

\textbf{Текст задания}

Вычислите:
\[
\left(\frac{3}{5}\right)^{-2}\cdot\frac{9}{25}.
\]

\textbf{Ваше решение}

В работе записано:
\[
\left(\frac{3}{5}\right)^{-2}\cdot\frac{9}{25}
=
\left(\frac{3}{5}\right)^{-2}\cdot\left(\frac{3}{5}\right)^2
=
\frac{3}{5}
=
0{,}6.
\]

\textbf{Ваш ответ}

\[
0{,}6.
\]

\textcolor{red}{\textbf{Ошибка:}} неверно применено правило умножения степеней с одинаковым основанием.

\textbf{Где именно возникает ошибка}

Последний правильный шаг:
\[
\left(\frac{3}{5}\right)^{-2}\cdot\frac{9}{25}
=
\left(\frac{3}{5}\right)^{-2}\cdot\left(\frac{3}{5}\right)^2.
\]

Здесь преобразование выполнено верно, потому что
\[
\frac{9}{25}=\frac{3^2}{5^2}=\left(\frac{3}{5}\right)^2.
\]

Первый неправильный шаг:
\[
\left(\frac{3}{5}\right)^{-2}\cdot\left(\frac{3}{5}\right)^2
=\frac{3}{5}.
\]

Именно это равенство неверно.

\textbf{Почему этот шаг неверен}

При умножении степеней с одинаковым основанием используется правило
\[
a^m\cdot a^n=a^{m+n},\qquad a\ne 0.
\]

В данном выражении основание в обеих степенях одинаковое --- это дробь $\frac{3}{5}$. Показатели равны $-2$ и $2$, поэтому их нужно именно сложить:
\[
-2+2=0.
\]

Следовательно,
\[
\left(\frac{3}{5}\right)^{-2}\cdot\left(\frac{3}{5}\right)^2
=
\left(\frac{3}{5}\right)^{-2+2}
=
\left(\frac{3}{5}\right)^0.
\]

Любое ненулевое число в нулевой степени равно единице. Дробь $\frac{3}{5}$ не равна нулю, поэтому
\[
\left(\frac{3}{5}\right)^0=1.
\]

Таким образом, переход к $\frac{3}{5}$ нарушает правило сложения показателей степеней. Показатели $-2$ и $2$ являются противоположными числами и полностью взаимно уничтожаются при сложении.

\textcolor{blue!60!black}{\textbf{Ключевая идея:}} если при умножении степени имеют одинаковое ненулевое основание, основание сохраняется, а показатели складываются. Здесь сумма показателей равна нулю, поэтому всё произведение равно единице.

\textbf{Исправление от последнего правильного шага}

Начинаем с уже верно полученного выражения:
\[
\left(\frac{3}{5}\right)^{-2}\cdot\left(\frac{3}{5}\right)^2.
\]

Складываем показатели:
\[
\left(\frac{3}{5}\right)^{-2+2}
=
\left(\frac{3}{5}\right)^0
=1.
\]

\textbf{Полное правильное решение}

Сначала представим дробь $\frac{9}{25}$ как степень с основанием $\frac{3}{5}$:
\[
\frac{9}{25}=\left(\frac{3}{5}\right)^2.
\]

Тогда исходное выражение принимает вид
\[
\left(\frac{3}{5}\right)^{-2}\cdot\left(\frac{3}{5}\right)^2.
\]

Основания одинаковые, поэтому при умножении складываем показатели:
\[
\left(\frac{3}{5}\right)^{-2+2}
=
\left(\frac{3}{5}\right)^0.
\]

Так как $\frac{3}{5}\ne 0$, получаем
\[
\left(\frac{3}{5}\right)^0=1.
\]

Следовательно,
\[
\boxed{1}.
\]

\textbf{Проверка}

Проверим результат другим способом, раскрыв отрицательную степень:
\[
\left(\frac{3}{5}\right)^{-2}
=
\left(\frac{5}{3}\right)^2
=
\frac{25}{9}.
\]

Тогда
\[
\frac{25}{9}\cdot\frac{9}{25}=1.
\]

Проверка даёт тот же результат.

\textcolor{teal!60!black}{\textbf{Итог:}} правильный ответ --- $1$. Ошибка возникла на шаге умножения степеней с одинаковым основанием: показатели $-2$ и $2$ нужно сложить, получив ноль.

\textbf{Задача на закрепление}

Вычислите:
\[
\left(\frac{2}{7}\right)^{-3}\cdot\frac{8}{343}.
\]

\vfill
\begin{center}
\small Лёвин Артём Александрович эксклюзивно для Матвея Горбачёва
\end{center}

\end{document}
